% GENERATED FILE. Edit canonical lessons, then run npm run generate:books.
% canonical-source-hash: fnv1a64:9537b3ea1291506c
% canonical-lessons: HI-C71-kyon, HI-C71-kyonki, HI-R71-reasons

\chapter{Why, and Because}
\label{ch:hi-why-because}
\hlchaptermodality{71}

\begin{chapteropening}
\textbf{By the end of this chapter:} \emph{I can ask why in Hindi and answer it with a whole reason clause after kyonki, including the one an exam message needs: why I cannot come.}

Complete HI-R71-reasons and take both turns unaided: ask kyon, then answer with the claim first and kyonki before the reason.
\end{chapteropening}

\section[kyoṁ]{\hi{क्यों} (\emph{kyoṁ}) — why}

\label{lesson:HI-C71-kyon}



\begin{quote}
Retrieve the two items before this one by name: \textbf{\hi{सकना}} and \textbf{\hi{चाहना}}. Then run the k- words the deixis lesson gave you — \textbf{\hi{क्या}}, \textbf{\hi{कहाँ}}, and \emph{kaun}. There is one more in that family, and it is the one a listener most often wants.
\end{quote}

\subsection*{You'll want to know: \hi{क्यों}}
\begin{quote}
\textbf{\hi{क्यों}} — \emph{kyoṁ} — \textbf{why}
\end{quote}

\begin{quote}
\textbf{\hi{आप} \hi{क्यों} \hi{नहीं} \hi{आ} \hi{सकते}?} — \emph{āp kyoṁ nahīṁ ā sakte?} — \textbf{Why can't you come?}
\end{quote}

\textbf{\hi{क्यों}} goes where any Hindi question word goes: in the sentence's ordinary slot, in front of the verb — \textbf{not} dragged to the front the way English drags \emph{why}. The rest of the sentence keeps the order it always had, verb last.

On its own, \textbf{\hi{क्यों}?} is a complete turn. It is how you ask somebody to explain what they have just said, and it costs one syllable.

The deixis lesson generalised the \textbf{\hi{य}- / \hi{व}- / \hi{क}-} system and listed the k- words. This is one of the ones it listed and did not teach. Its \textbf{\hi{क}-} is the same ancient stem as \textbf{\hi{क्या}} and \emph{kaun}, and it is Latin's \textbf{qu-} in \emph{quis} and English's \textbf{wh-} in \emph{who} — one Indo-European \textbf{k\textsuperscript{w}-}, three different spellings of the same sound.

\subsection*{Your turn}
Say these aloud:

\begin{itemize}
\raggedright
  \item \emph{kyoṁ?} on its own, as a whole turn
  \item \emph{āp kyoṁ nahīṁ ā sakte?}
  \item the four k- words together — \emph{kyā}, \emph{kaun}, \emph{kahāṁ}, \emph{kyoṁ}
\end{itemize}

\begin{tcolorbox}[breakable,colback=teal!4,colframe=teal!35!black,title={Before you move on}]
Where does \textbf{\hi{क्यों}} stand in the sentence? (\textbf{In its ordinary slot, in front of the verb — not fronted}.) Which family does it belong to? (\textbf{The k- question words}.)

Sources: \href{https://en.wiktionary.org/wiki/\%E0\%A4\%95\%E0\%A5\%8D\%E0\%A4\%AF\%E0\%A5\%8B\%E0\%A4\%82}{Wiktionary: \hi{क्यों}}.
\end{tcolorbox}
\section[kyoṁki]{\hi{क्योंकि} (\emph{kyoṁki}) — because}

\label{lesson:HI-C71-kyonki}



\begin{quote}
Ask \textbf{\hi{क्यों}?} out loud. Then notice that you have no way at all to answer it. Now say \textbf{\hi{कि}} — the word from the last chapter that hangs a clause on a verb. The answer is those two words stuck together.
\end{quote}

\begin{grammarlens}[title={Grammar Lens: \hi{क्यों} + \hi{कि}}]
\begin{quote}
\textbf{\hi{मैं} \hi{कल} \hi{नहीं} \hi{आ} \hi{सकता} \hi{क्योंकि} \hi{मैं} \hi{काम} \hi{करता} \hi{हूँ}.} — \emph{maiṁ kal nahīṁ ā saktā kyoṁki maiṁ kām kartā hūṁ.} — \textbf{I can't come tomorrow because I work.}
\end{quote}

Read the word: \textbf{\hi{क्योंकि}} is \textbf{\hi{क्यों}} + \textbf{\hi{कि}}, \emph{why-that}. Nothing is hidden in it. Hindi answers \emph{why} by saying \emph{why-that} and then a whole ordinary sentence, verb last, exactly as after plain \textbf{\hi{कि}}.

The order is fixed and it is the useful way round: \textbf{the claim first, the reason second}. English can lead with \emph{because}; the ordinary Hindi sentence does not.

That example is worth memorising whole, because it is the shape of an entire exam task: a written message that says you cannot meet and why. It needs three things this book has now given you — the negation, \textbf{\hi{सकना}}, and \textbf{\hi{क्योंकि}} — and it needed all three.

Look at what the word is made of one more time. \textbf{\hi{क्यों}} comes down from Sanskrit; \textbf{\hi{कि}} was borrowed from Persian. Hindi welded them into a single conjunction and nobody notices the seam. That is this language in one word.
\end{grammarlens}

\subsection*{Your turn}
Say these aloud:

\begin{itemize}
\raggedright
  \item the two pieces separately — \emph{kyoṁ}, \emph{ki} — then welded, \emph{kyoṁki}
  \item \emph{maiṁ kal nahīṁ ā saktā kyoṁki maiṁ kām kartā hūṁ}
  \item which half of that sentence comes first — \textbf{the claim}
\end{itemize}

\begin{tcolorbox}[breakable,colback=teal!4,colframe=teal!35!black,title={Before you move on}]
What two words is \textbf{\hi{क्योंकि}} made of? (\textbf{\hi{क्यों} + \hi{कि}}.) Which half of a \textbf{\hi{क्योंकि}} sentence comes first? (\textbf{The claim; the reason follows}.)

Sources: \href{https://en.wiktionary.org/wiki/\%E0\%A4\%95\%E0\%A5\%8D\%E0\%A4\%AF\%E0\%A5\%8B\%E0\%A4\%82\%E0\%A4\%95\%E0\%A4\%BF}{Wiktionary: \hi{क्योंकि}}.
\end{tcolorbox}
\section[kyoṁ? — kyoṁki …]{Two turns: why, and because}

\label{lesson:HI-R71-reasons}



\begin{quote}
\textbf{\hi{क्योंकि}} is two words you already had. Say them apart, then together, then say where the verb of each half sits.

Reach back three chapters while you are here: \textbf{\hi{या}} offers a choice and \textbf{\hi{लेकिन}} contradicts. Neither can give a reason, which is what this chapter added.
\end{quote}

\subsection*{Your turn}
Say these aloud:

\begin{itemize}
\raggedright
  \item \emph{kyoṁ?} — then \emph{kyoṁki …}
  \item \emph{maiṁ kal nahīṁ ā saktā kyoṁki maiṁ kām kartā hūṁ}
  \item \emph{maiṁ ānā chāhtā hūṁ lekin maiṁ nahīṁ ā saktā} — the joiner and the two modal verbs in one sentence
\end{itemize}

\begin{tcolorbox}[breakable,colback=teal!4,colframe=teal!35!black,title={Before you move on}]
Which word asks for a reason, and which gives one? (\textbf{\hi{क्यों}} asks, \textbf{\hi{क्योंकि}} gives.) What is the second one made of? (\textbf{The first one, plus \hi{कि}}.) Then stop.
\end{tcolorbox}
